\problema{Course Schedule II}{210}{src/01-grafos/210-course-schedule-ii.cpp}{M}

Hay que cursar \texttt{numCourses} materias numeradas de \texttt{0} a
\texttt{numCourses-1}. Nos dan una lista de prerrequisitos
\texttt{prerequisites[i] = [a, b]}, que significa que para tomar el curso
\texttt{a} primero hay que haber tomado el curso \texttt{b}. Queremos
devolver \emph{un} orden válido para cursar todas las materias, o un vector
vacío si es imposible (existe un ciclo de dependencias).

\paragraph{Intuición.} Es la continuación natural del problema 207 (Course
Schedule): ahí solo preguntaban si era \emph{posible} terminar todos los
cursos; aquí además hay que construir el orden concreto. La idea sigue
siendo la misma: un curso solo se puede tomar cuando ya se tomaron todos sus
prerrequisitos, así que vamos "liberando" cursos conforme se cumplen sus
dependencias, con el algoritmo de Kahn (BFS sobre el grado de entrada).

\begin{center}
\ttfamily
\begin{tabular}{ccc}
 & 0 & \\
1 & & 2\\
 & 3 & \\
\end{tabular}
\end{center}

\noindent Con \texttt{prerequisites = [[1,0],[2,0],[3,1],[3,2]]}: para tomar
\texttt{1} y \texttt{2} hace falta \texttt{0}; para tomar \texttt{3} hacen
falta \texttt{1} y \texttt{2}. Un orden válido es \texttt{[0, 1, 2, 3]} (o
\texttt{[0, 2, 1, 3]}): el orden topológico no es único cuando hay ramas
independientes.

\paragraph{Solución.} Construimos la lista de adyacencia
\texttt{curso -> [dependientes]} y el grado de entrada de cada curso
(cuántos prerrequisitos le faltan). Metemos en una cola todos los cursos con
grado de entrada \texttt{0}; al sacar un curso de la cola lo agregamos al
resultado y restamos \texttt{1} al grado de entrada de sus dependientes, y el
que llegue a \texttt{0} entra a la cola. Si al final el resultado no incluye
a los \texttt{numCourses} cursos, hay un ciclo y devolvemos un vector vacío.

\lstinputlisting{src/01-grafos/210-course-schedule-ii.cpp}

\complejidad{$O(V+E)$}{$O(V+E)$}

\paragraph{Notas de entrevista (Amazon).} Como el orden topológico casi
nunca es único, en la entrevista aclara que cualquier orden que respete las
dependencias es una respuesta válida; si te piden verificar tu propia
solución, basta con checar que cada prerrequisito aparezca antes que su
curso. Casos borde a mencionar: sin prerrequisitos (cualquier permutación
sirve), cadena lineal (orden único), y ciclo (respuesta vacía, igual que
"false" en el 207).
